\documentclass[11pt,a4paper]{article}
\usepackage[margin=19mm]{geometry}
\usepackage{fontspec}
\setmainfont{Latin Modern Roman}
\setmonofont{Latin Modern Mono}
\newfontfamily\taafont{DejaVu Sans}
\usepackage{amsmath,amssymb,amsthm,microtype}
\usepackage[hidelinks]{hyperref}
\newtheorem{definition}{Definition}
\newtheorem{theorem}{Theorem}
\newtheorem{proposition}{Proposition}
\newcommand{\TAA}{\mathsf{TAA}}
\hypersetup{pdftitle={Exact State Contraction: Definition and Notation},pdfauthor={lmtlssss}}
\title{{\taafont ᑕ.}\\Exact State Contraction\\\large A definition, preservation theorem, and notation proposal}
\author{lmtlssss}
\date{8 November 2026 \quad Version 1.0}
\begin{document}
\maketitle
\begin{abstract}
We define exact state contraction relative to a declared observation contract.
The retained transition and observation maps commute with the full system.
This identity preserves every finite observation trace by induction.
An affine interval extension carries a quadratic accumulated cost through the
same associative composition, including cross-terms. We distinguish output
preservation from witness-count preservation and specify the printed notation
{\taafont ᑕ.}, spoken TAA. This paper names a contract and gives its elementary
preservation theorem; it does not claim priority for quotient systems,
bisimulation, or affine composition.
\end{abstract}

\section{The contract}
Let $M=(S,U,F,Y,G)$ be an input-driven system. A state $s\in S$ and an
input $u\in U$ produce the next state $F(s,u)$ and an observation $G(s,u)$.
The observation identifies the property to preserve.
\begin{definition}[Exact state contraction]
An exact state contraction of $M$ for $G$ consists of a set $Q$, a
surjection $\pi:S\rightarrow Q$, and maps
$\bar F:Q\times U\rightarrow Q$ and $\bar G:Q\times U\rightarrow Y$ such that
\begin{equation}\label{eq:contract}
 \pi(F(s,u))=\bar F(\pi(s),u),\qquad
 G(s,u)=\bar G(\pi(s),u)
 \quad (s\in S,u\in U).
\end{equation}
We denote the tuple by $\TAA_G(M)=(Q,\pi,\bar F,\bar G)$.
Its printed operator name is {\taafont ᑕ.}, pronounced TAA.
\end{definition}

The subscript $G$ is necessary. A contraction can preserve one observation
while losing information required for another.

\begin{theorem}[Trace preservation]
For each initial state $s_0$ and finite input word, the system $M$ and
its contraction give equal observations at every step when the retained
initial state is $\pi(s_0)$.
\end{theorem}
\begin{proof}
Assume $q_k=\pi(s_k)$. The first identity in (\ref{eq:contract}) gives
$q_{k+1}=\bar F(q_k,u_k)=\pi(F(s_k,u_k))=\pi(s_{k+1})$.
The second identity gives
$\bar G(q_k,u_k)=G(s_k,u_k)$.
Since $q_0=\pi(s_0)$, induction proves both statements for each step.
\end{proof}

For relations with multiple witnesses, preservation of the set of output
values does not prove preservation of witness multiplicity. A claim of
unique witness requires a bijection of the applicable witness fibers or an
independent exact count. A claim of unique output requires equality of the
applicable output sets.

\section{Retained affine cost}
An interval $A$ carries the state transition
$T_A(x)=F_Ax+b_A$ and the cost
$J_A(x)=x^\mathsf{T}Q_Ax+q_A^\mathsf{T}x+d_A$, with $Q_A$ symmetric.
For interval $B$ after $A$, define
$T_{BA}=T_B\circ T_A$ and $J_{BA}=J_A+J_B\circ T_A$.
\begin{proposition}[Quadratic energy composition]
The composed coefficients are
\begin{align*}
F_{BA}&=F_BF_A,&b_{BA}&=F_Bb_A+b_B,\\
Q_{BA}&=Q_A+F_A^\mathsf{T}Q_BF_A,\\
q_{BA}&=q_A+F_A^\mathsf{T}(q_B+2Q_Bb_A),\\
d_{BA}&=d_A+d_B+q_B^\mathsf{T}b_A+b_A^\mathsf{T}Q_Bb_A.
\end{align*}
This composition is associative and gives the exact chronological cost
for each input state.
\end{proposition}
\begin{proof}
Substitute $F_Ax+b_A$ into $J_B$ and collect quadratic,
linear, and constant coefficients. The displayed formulas result.
For a third interval $C$, both groupings give the transition
$T_C\circ T_B\circ T_A$ and cost
$J_A+J_B\circ T_A+J_C\circ T_B\circ T_A$.
Therefore both groupings agree as functions.
\end{proof}

\paragraph{Example.}
Let $T_A(x)=2x+1$, $J_A(x)=x^2+1$,
$T_B(y)=3y-2$, and $J_B(y)=2y^2+3y+4$.
Then $T_{BA}(x)=6x+1$ and
$J_{BA}(x)=9x^2+14x+10$.
At $x=2$, the direct costs are 5 and 69, giving 74.
The retained polynomial also gives 74. The companion verifier checks
the exact rational identities and a family of associative compositions.

\section{Certificates and meaning}
A {\taafont ᑕ.} certificate states the artifact digest, claim,
admissible input domain, observation contract, checker version, verdict,
and evidence digest. The signature authenticates the issuer's statement;
the independent checker establishes whether the evidence supports the
declared claim. A transparency receipt records registration under one
service's policy. None of these substitutes for the others.

The word \emph{exact} is relative to the declared observation and domain.
The definition gives no claim that arbitrary systems admit a small
retained state or that arbitrary program properties are efficiently
decidable.

\section{Symbol and prior results}
The symbol {\taafont ᑕ.} consists of U+1455 CANADIAN SYLLABICS TA and
U+002E FULL STOP. The existing glyph has a linguistic identity in the
Unified Canadian Aboriginal Syllabics block. This paper assigns a
mathematical use to those characters; it does not claim to create them.
The software identifier is \texttt{taa}, with the versioned name
\texttt{taa.exact-state-contraction.v1}.

The author's earlier work proved retained affine-cost identities,
minimal linear recovery ranks, and exact coefficient certificates for
specific contracts \cite{affine}. Established quotient structures,
bisimulation, automata minimization, and affine composition remain
the mathematical background of this proposal.

\begin{thebibliography}{9}
\bibitem{affine} lmtlssss, \emph{Affine Cycles: Exact Cut Defects and Minimal
Retained State}, Zenodo (2026), DOI: \href{https://doi.org/10.5281/zenodo.23931722}{10.5281/zenodo.23931722}.
\bibitem{unicode} Unicode Consortium, \emph{Unified Canadian Aboriginal Syllabics}, Unicode Character Database, U+1455.
\bibitem{scitt} IETF, \emph{An Architecture for Trustworthy and Transparent Digital Supply Chains}, RFC 9943 (2026).
\end{thebibliography}
\end{document}
